\lyseis{Β1.13}{Θέσεις ευθείας και κύκλου}
\sol{1} α) εξωτερική \quad β) εφαπτομένη \quad γ) τέμνουσα \quad δ) τέμνουσα (διάμετρος) \quad ε) εξωτερική.
\sol{2} $0$, $1$ και $2$ κοινά σημεία αντίστοιχα.
\sol{3} Ελεύθερη κατασκευή. Για την εφαπτομένη το $d$ βγαίνει $3$ cm.
\sol{4} Φέρνω την ακτίνα $\mathrm{ΟΜ}$ και την κάθετη στην $\mathrm{ΟΜ}$ στο $\mathrm{Μ}$.
\sol{5} $0\le d<5$.
\sol{6} α) Σωστό. β) Λάθος: δύο. γ) Σωστό. δ) Λάθος: είναι εξωτερική.
\sol{7} Όχι: $d=\rho$, η ευθεία εφάπτεται και ακουμπά μόνο σε ένα σημείο.
\sol{8} $1{,}8$ μ.: τέμνει. $2{,}5$ μ.: εφάπτεται. $3{,}4$ μ.: δεν το ακουμπά.
\sol{9} $33$ cm (η ακτίνα). Ένα κοινό σημείο.
\sol{10} Στο σημείο επαφής του φωτός με τον τοίχο, ένα μόνο σημείο.
\sol{11} Ο δρόμος εφάπτεται στη ρόδα: υπάρχει μόνο ένα σημείο επαφής και το κέντρο μένει σε σταθερή απόσταση ίση με την ακτίνα.
\sol{12} Η εφαπτομένη έχει με τον κύκλο ένα μόνο κοινό σημείο, και όλα τα άλλα σημεία της απέχουν από το κέντρο περισσότερο από $\rho$. Άρα το $\mathrm{Μ}$ είναι το πλησιέστερο, και το πλησιέστερο σημείο είναι εκεί που πέφτει η κάθετη από το κέντρο.
\sol{13} Όχι: από δύο σημεία περνά μία ευθεία, και μια ευθεία έχει με κύκλο το πολύ δύο κοινά σημεία.
\sol{14} Από σημείο του κύκλου περνά μία εφαπτομένη. Από σημείο εκτός κύκλου περνούν δύο.
\sol{15} Η απόσταση είναι $2\rho$.
\sol{16} Ελεύθερη απάντηση.
\sol{17} Ελεύθερη κατασκευή.
